\section{The window correction}\label{sec:window}

\paragraph{The truncation artifact.} In the CFT every descendant spacing $E(\partial^n\mathcal O)-E(\mathcal O)=n$
is exact. The raw truncation violates these spacings by amounts that do not decay with $N$. In $d=1$ the level-1
violations are $\HTq{V1}$ at $\Ds\le0.1$ over $N=16$--28 (Table~\ref{tab:d1-ident} shows $N=28$). The mechanism is that a bath state of level $\ell$ can carry
vacuum fluctuations of the bath only up to level $N-\ell$, whereas the vacuum keeps all fluctuations up to $N$. The
same mechanism spoils the cancellation of disconnected vacuum diagrams in $\phi^4$ theory in $2+1$ dimensions with an
energy cutoff, where the effect does not decouple as the cutoff is removed; there it was addressed with perturbative
state-dependent counterterms~\cite{Elias-Miro:2020qwz}. The missing vacuum window is the truncation artifact; in
the perturbative truncation it is analysed in \texttt{HT\_Pert/}.

\paragraph{The correction.} We add to every bath state of level $\ell$, in all sectors, the shift
\begin{equation}
 D(\ell)=E_v(N)-E_v(N-\ell),\qquad D(0)=0,
 \label{eq:window}
\end{equation}
where $E_v(M)$ is the vacuum energy at bath cutoff $M$ of the same Hamiltonian, with the same $\Ds$, couplings, zero
mode and basis. The definition of $E_v$ has three features:
\begin{itemize}
\item \emph{Zero mode included.} Vacua of the bath alone miss the zero-mode tadpole $\frac g2\langle q^2\rangle$,
 which cancels the $-\frac g2$ in $A$, and remove only part of the violation.
\item \emph{Consistent.} The production mode is \emph{consistent}: $E_v(M)$ is the lowest eigenvalue of the
 \emph{corrected} Hamiltonian at cutoff $M$,
 \begin{equation}
 E_v(M)=\lambda_{\min}\big[H(M)+D_M\big],\qquad D_M(\ell)=E_v(M)-E_v(M-\ell),
 \label{eq:Ev}
 \end{equation}
 built level by level from the corrected lower cutoffs.
\item \emph{Odd cutoffs.} $E_v$ of an odd $M$ equals $E_v(M-1)$, since only even levels exist in the sector.
\end{itemize}
The \emph{literal} mode uses the uncorrected $\lambda_{\min}H(M)$ and serves as a diagnostic. Energies are measured
from the lowest even eigenvalue of the zero mode alone, so $E_v(0)=0$.

\paragraph{Precedent.} $D(\ell)$ is the counterpart, with a level cutoff, of the state-dependent counterterms that
Ref.~\cite{Elias-Miro:2020qwz} adds with an energy cutoff to restore the missing states of disconnected bubbles.
\begin{itemize}
\item At second order in the interaction, $D(\ell)$ corresponds to their single-bubble counterterm (Patch~I, their
 Eq.~(5.38)). Their denominators contain the unperturbed energy of the level being computed; ours contain the bubble
 energy alone.
\item They add such terms order by order, including towers of repeated bubbles (Patch~II, up to five bubbles), as a
 pragmatic fix. $D(\ell)$ is a nonperturbative analogue: it takes all disconnected bubbles at once from the vacuum
 energies of the same truncated Hamiltonian at lower cutoffs. The consistent definition, Eq.~\eqref{eq:Ev}, also
 corrects the states inside those vacuum energies, which is the role of their Patch~II.
\end{itemize}

\paragraph{Fixed point.} At each $M$, Eq.~\eqref{eq:Ev} is a scalar fixed point $s=\lambda_{\min}[H(M)+D_M(s)]$,
because $s=E_v(M)$ enters every $D_M(\ell>0)$. We write $F(s)=\lambda_{\min}-s$. Then $F'(s)=-p_0$, with $p_0$ the
weight of the vacuum vector on bath level 0. $F$ is concave, so Newton's method converges monotonically after at most
one step, and the fixed points are solved to $10^{-12}$. At $M=N$, $s$ is the ground energy of the main problem.

\paragraph{Derivatives.} For the tuning Jacobian, Hellmann--Feynman on the stored vacuum vectors gives
\begin{equation}
 \partial E_v(M)=\frac{\langle\partial H\rangle-\sum_{\ell>0}w_\ell\,\partial E_v(M-\ell)}{p_0},
 \label{eq:dEv}
\end{equation}
with $w_\ell$ the weight of the vacuum vector on bath level $\ell$. The package tests compare it with finite differences.

\paragraph{What it removes and what it misses.}
\begin{itemize}
\item In $d=1$ the correction reduces the level-1 violation of $\sigma$ by a factor $\HTq{V2}$ at $N=28$,
 $\Ds\le0.15$.
\item At $\eps\ge0.6$ ($\Ds\le0.1$), non-decaying residuals remain:
 \begin{itemize}
 \item about $\HTq{V3}$ for $\sigma$ and $\phi^3$;
 \item up to $\HTq{V4}$ for $\varepsilon$ at $\Ds\le0.05$ (Table~\ref{tab:d1-ident}).
 \end{itemize}
\item The correction restores only disconnected vacuum fluctuations. At finite coupling it therefore models the
 truncation artifact rather than removing it completely, and the residuals show what it misses.
\item Expanded in $\eps$, the corrected $d=1$ model changes $\DE$ only at $O(\eps^4)$: in the blind solver the
 difference to the raw $\DE$ scales as $\eps^{4}$ (Section~\ref{sec:checks}). $C_2$ and $C_3$ are therefore unaffected.
\item The perturbative windows of \texttt{HT\_Pert/} are built in the same spirit. In $d=1$ the finite-coupling
 window has not been compared with them order by order, so whether its $O(\eps^4)$ term is the one the continuum
 requires is not known.
\item In $d=2,3$ the consistent window is the same construction (Sections~\ref{d2:sec}, \ref{d3:sec}). In $d=2$ its
 $\eps$ expansion reproduces the consistent-window scheme of \texttt{HT\_Pert/}, whose fourth-order coefficient
 grows like $\ln N$; the model is therefore correct only through third order. In $d=3$ its coefficients $C_n(N)$
 come from \texttt{HT\_Pert/}: $C_2$ and $C_3$ approach their continuum values, while the fourth-order coefficient at
 $N=20$ is 60\% of $C_4(3)$ and still growing slowly, a finite-$N$ statement.
\end{itemize}
We have not constructed an effective Hamiltonian beyond the window~\cite{Cohen:2021erm,EliasMiro:2022pua}.
